\documentclass[11pt,letterpaper]{article}

% ── Packages ──────────────────────────────────────────────────────────────
\usepackage[utf8]{inputenc}
\usepackage[T1]{fontenc}
\usepackage{lmodern}
\usepackage[margin=1in]{geometry}
\usepackage{amsmath,amssymb,amsthm}
\usepackage{booktabs}
\usepackage{enumitem}
\usepackage{graphicx}
\usepackage{hyperref}
\usepackage{xcolor}
\usepackage{listings}
\usepackage{fancyhdr}
\usepackage{tcolorbox}

% ── Colors ────────────────────────────────────────────────────────────────
\definecolor{codeblue}{RGB}{46,134,171}
\definecolor{codegray}{RGB}{100,100,100}
\definecolor{codegreen}{RGB}{68,187,164}
\definecolor{codebg}{RGB}{248,248,248}
\definecolor{warnred}{RGB}{232,72,85}

% ── Listings style for C++ ────────────────────────────────────────────────
\lstdefinestyle{cpp}{
    language=C++,
    basicstyle=\ttfamily\small,
    keywordstyle=\color{codeblue}\bfseries,
    commentstyle=\color{codegray}\itshape,
    stringstyle=\color{codegreen},
    backgroundcolor=\color{codebg},
    frame=single,
    rulecolor=\color{codegray!30},
    numbers=left,
    numberstyle=\tiny\color{codegray},
    numbersep=8pt,
    breaklines=true,
    tabsize=4,
    showstringspaces=false,
    columns=flexible,
    morekeywords={Eigen,MatrixXd,VectorXd,SparseMatrix,Triplet,
                  SparseLU,BiCGSTAB,IncompleteLUT,ComputeThinU,
                  ComputeThinV,Success}
}
\lstset{style=cpp}

% ── Theorem-like environments ─────────────────────────────────────────────
\theoremstyle{definition}
\newtheorem{definition}{Definition}[section]
\newtheorem{example}{Example}[section]
\theoremstyle{remark}
\newtheorem{remark}{Remark}[section]

% ── Warning box ───────────────────────────────────────────────────────────
\newtcolorbox{warningbox}{
    colback=warnred!5,
    colframe=warnred!70,
    fonttitle=\bfseries,
    title={Warning}
}

% ── Tip box ───────────────────────────────────────────────────────────────
\newtcolorbox{tipbox}{
    colback=codegreen!5,
    colframe=codegreen!70,
    fonttitle=\bfseries,
    title={Engineer's Tip}
}

% ── Header / Footer ──────────────────────────────────────────────────────
\pagestyle{fancy}
\fancyhf{}
\lhead{\textsc{Practical NLA}}
\rhead{\thepage}
\renewcommand{\headrulewidth}{0.4pt}

% ── Title ─────────────────────────────────────────────────────────────────
\title{\textbf{Numerical Linear Algebra for Engineers}\\[6pt]
       \large The Combat Medic Guide:\\
       Physics $\rightarrow$ Code $\rightarrow$ Result}
\author{}
\date{}

\begin{document}
\maketitle
\thispagestyle{empty}

\begin{abstract}
This is not a math textbook. This is a field manual. We cut the proofs, keep the
tools, and show you how to solve real linear systems in C++ using Eigen. Every
section follows the same arc: \textbf{here is the physics}, \textbf{here is why
the naive approach fails}, and \textbf{here is the code that works}.
\end{abstract}

\tableofcontents
\newpage

% ══════════════════════════════════════════════════════════════════════════
% §1  The Concrete Problem: The 1-D Heat Rod
% ══════════════════════════════════════════════════════════════════════════
\section{The Concrete Problem: The 1-D Heat Rod}

\subsection{The Physics}

We have a metal rod of length $L$.  Both ends are held at $0^\circ$C.  A heat
source $f(x)$ is applied along the rod.  We want the \textbf{steady-state}
temperature $u(x)$, which satisfies the Poisson equation
\[
    -u''(x) = f(x), \qquad u(0) = u(L) = 0.
\]
Discretise the rod into $n$ interior points with spacing $h = L/(n+1)$.
Replacing $u''$ with the second-order central difference gives
\[
    \frac{-u_{i-1} + 2u_i - u_{i+1}}{h^2} = f_i,
    \qquad i = 1,\dots,n.
\]
This is a linear system $Au = h^2 f$ with the \textbf{tridiagonal} matrix
\begin{equation}\label{eq:tridiag}
    A = \begin{bmatrix}
         2 & -1 &        &        \\
        -1 &  2 & -1     &        \\
           & \ddots & \ddots & \ddots \\
           &        & -1     &  2
    \end{bmatrix}
    \in \mathbb{R}^{n\times n}.
\end{equation}

% ── Case A ────────────────────────────────────────────────────────────────
\subsection{Case A: The $3\times 3$ Problem (By Hand)}

Set $n=3$ and $f_i=1$ for all~$i$.  Then $b = h^2[1,1,1]^T$ and
\[
    A = \begin{bmatrix} 2 & -1 & 0 \\ -1 & 2 & -1 \\ 0 & -1 & 2 \end{bmatrix}.
\]
You can solve this by Gaussian elimination in a few lines.
The exact solution is $u = \tfrac{1}{4}[1, \tfrac{3}{2}, 1]^T$ (after scaling
by $h^2$).  Computing $A^{-1}$ directly already requires tracking nine entries
and two row-reduction passes---even a single rounding mistake
in one pivot cascades through every subsequent row.

\begin{remark}
For a \emph{tridiagonal} system, Gaussian elimination is $O(n)$.
General $LU$ is $O(n^3)$.  Structure \emph{always} matters.
\end{remark}

The following C++ snippet solves Case~A with Eigen:

\lstinputlisting[caption={Case A: $3\times 3$ heat rod (dense).},
                 label={lst:heat3}]
                {code/heat_rod_3x3.cpp}

% ── Case B ────────────────────────────────────────────────────────────────
\subsection{Case B: The $1{,}000{,}000 \times 1{,}000{,}000$ Problem}

Now the rod has $n = 10^6$ interior points.

\begin{itemize}[nosep]
    \item A \textbf{dense} $10^6 \times 10^6$ matrix of doubles needs
          $10^{12} \times 8\;\text{B} = 8\;\text{TB}$ of RAM.
    \item The tridiagonal $A$ has only $\approx 3n$ non-zeros---about
          $24\;\text{MB}$.
\end{itemize}

\begin{warningbox}
If you try to allocate a dense $10^6 \times 10^6$ matrix,
your program will be killed by the OS before it prints ``Hello.''
\end{warningbox}

This single observation motivates every topic that follows:
\emph{sparse storage} (\S\ref{sec:sparse}),
\emph{iterative solvers} (\S\ref{sec:gmres}),
and \emph{hardware awareness} (\S\ref{sec:hardware}).

\lstinputlisting[caption={Case B: $10^6$ heat rod (sparse, direct solve).},
                 label={lst:heat1M}]
                {code/heat_rod_1M.cpp}

% ══════════════════════════════════════════════════════════════════════════
% §2  Hardware Reality: Why Your O(n³) is Slow
% ══════════════════════════════════════════════════════════════════════════
\section{Hardware Reality: Why Your \texorpdfstring{$O(n^3)$}{O(n³)} is Slow}
\label{sec:hardware}

In C++, \textbf{memory is the bottleneck, not arithmetic}.
A modern CPU can execute $\sim\!10^{10}$ floating-point operations per second,
but fetching a cache line from main RAM takes $\sim\!100$ ns---long enough
for $\sim\!200$ multiplies to have completed.

\subsection{Cache Locality}

CPUs load data in \emph{cache lines} (typically 64 bytes = 8 doubles).
If your next access is a neighbour in memory, it is already in cache
(\textbf{cache hit}).  If you jump across rows or columns, the CPU stalls
while fetching a new line (\textbf{cache miss}).

\begin{definition}[Row-major vs.\ Column-major]
In \textbf{row-major} layout (C/C++ default), the element
$A_{i,j}$ is stored at offset $i \cdot n + j$.
Incrementing $j$ visits contiguous memory.

In \textbf{column-major} layout (Fortran, MATLAB, \textbf{Eigen default}),
the element $A_{i,j}$ is stored at offset $j \cdot m + i$.
Incrementing $i$ visits contiguous memory.
\end{definition}

\begin{warningbox}
Eigen stores matrices in \textbf{column-major} order by default.
If you write a row-traversal loop \texttt{for(i) for(j) A(i,j)},
every inner-loop access jumps by $m$ doubles---a cache miss every time.
Swap the loop order, or use \texttt{Eigen::RowMajor}.
\end{warningbox}

The following benchmark makes the difference concrete:

\lstinputlisting[caption={Cache locality: row-major vs.\ column-major traversal.},
                 label={lst:cache}]
                {code/cache_locality.cpp}

\subsection{SIMD: Single Instruction, Multiple Data}

Modern x86 CPUs have \textbf{AVX-2} (4 doubles/cycle) or \textbf{AVX-512}
(8 doubles/cycle) vector units.  When your data is \emph{contiguous and
aligned}, the compiler auto-vectorises your loop and processes 4--8 elements
per instruction.  When data is scattered, the vector unit idles.

\begin{tipbox}
You almost never write SIMD intrinsics by hand.
The rule is simpler: \textbf{keep your data contiguous}, and the
compiler (with \texttt{-O2 -march=native}) will vectorise for you.
Libraries like Eigen and MKL are already written to exploit this.
\end{tipbox}

\subsection{The Takeaway}

\begin{center}
\renewcommand{\arraystretch}{1.3}
\begin{tabular}{@{}lcc@{}}
\toprule
\textbf{Access pattern} & \textbf{Cache behaviour} & \textbf{Relative speed} \\
\midrule
Contiguous (correct order) & Hits    & $1\times$ \\
Stride-$n$ (wrong order)   & Misses  & $3$--$10\times$ slower \\
Random access              & Misses  & $10$--$50\times$ slower \\
\bottomrule
\end{tabular}
\end{center}

% ══════════════════════════════════════════════════════════════════════════
% §3  The Library Layer: BLAS & Eigen
% ══════════════════════════════════════════════════════════════════════════
\section{The Library Layer: BLAS \& Eigen}
\label{sec:blas}

\subsection{Why Libraries Beat Hand-Written Loops}

The \textbf{Basic Linear Algebra Subprograms (BLAS)} specification defines
three levels of operations:

\begin{center}
\renewcommand{\arraystretch}{1.3}
\begin{tabular}{@{}clcc@{}}
\toprule
\textbf{Level} & \textbf{Operation} & \textbf{FLOPs} & \textbf{Data} \\
\midrule
1 & $\alpha x + y$ \quad (AXPY)       & $O(n)$   & $O(n)$ \\
2 & $Ax + y$ \quad (GEMV)             & $O(n^2)$ & $O(n^2)$ \\
3 & $\alpha AB + \beta C$ \quad (GEMM) & $O(n^3)$ & $O(n^2)$ \\
\bottomrule
\end{tabular}
\end{center}

\begin{tipbox}
Level~3 (GEMM) is the ``Gold Standard'' because its arithmetic intensity
is $O(n^3)/O(n^2) = O(n)$: the CPU reuses each cache line many times before
evicting it.  This is why matrix--matrix multiply runs close to peak FLOP/s,
while matrix--vector multiply is memory-bound.
\end{tipbox}

Optimised BLAS implementations---\textbf{Intel MKL}, \textbf{OpenBLAS},
\textbf{BLIS}---apply cache-blocking, SIMD, and multi-threading
automatically.  A na\"ive triple loop achieves $<5\%$ of peak;
a tuned GEMM reaches $>90\%$.

\subsection{The SpMV Speed Trap}

Engineers often assume that because they are using a ``fast'' library like
Eigen or MKL, their sparse solver will hit peak FLOP/s.
\textbf{It will not.}

When you solve a sparse system iteratively (e.g.\ GMRES on the
$10^6 \times 10^6$ heat rod), the inner loop is a \emph{Sparse
Matrix--Vector multiply} (SpMV)---a Level~2 operation.  Its arithmetic
intensity is $O(\text{nnz})/O(\text{nnz}) = O(1)$: each non-zero is
loaded from RAM, multiplied once, and never reused.

\begin{warningbox}
SpMV is \textbf{memory-bandwidth bound}, not compute-bound.
No amount of SIMD or MKL magic can fix a slow RAM bus if you are
just streaming a million doubles once per iteration.  On a typical
desktop with 40~GB/s bandwidth, SpMV on a tridiagonal $10^6$ system
peaks at $\sim\!3$~GFLOP/s---less than $5\%$ of the CPU's
theoretical throughput.
\end{warningbox}

This is why \emph{preconditioning} (\S\ref{sec:gmres}) matters so much:
reducing the iteration count by $100\times$ is worth far more than any
micro-optimisation of the SpMV kernel itself.

\subsection{Eigen as a BLAS Wrapper}

When you write
\begin{lstlisting}
Eigen::MatrixXd C = A * B;
\end{lstlisting}
Eigen's \emph{expression templates} detect this is a GEMM and dispatch to
whichever BLAS backend is linked at compile time.  The cascade is:

\begin{enumerate}[nosep]
    \item \textbf{MKL linked?} $\rightarrow$ calls \texttt{cblas\_dgemm}
          (fastest on Intel hardware).
    \item \textbf{OpenBLAS linked?} $\rightarrow$ calls its GEMM
          (good cross-platform default).
    \item \textbf{Neither?} $\rightarrow$ Eigen's built-in blocked GEMM
          (decent, but slower than the above).
\end{enumerate}

\begin{warningbox}
If you install Eigen but never link a BLAS library, you are leaving
a $2$--$5\times$ speedup on the table for every dense multiply.
\end{warningbox}

\subsection{Linking MKL in Practice}

With CMake and Intel oneAPI installed:
\begin{lstlisting}[language=bash,morekeywords={find_package,target_link_libraries}]
# CMakeLists.txt
find_package(MKL REQUIRED)
target_link_libraries(my_app PRIVATE MKL::MKL)
add_definitions(-DEIGEN_USE_MKL_ALL)
\end{lstlisting}

The following snippet benchmarks na\"ive loops against Eigen's GEMM:

\lstinputlisting[caption={BLAS-level GEMM: na\"ive loop vs.\ Eigen.},
                 label={lst:gemm}]
                {code/gemm_benchmark.cpp}

% ══════════════════════════════════════════════════════════════════════════
% §4  The Direct Hammer: SVD
% ══════════════════════════════════════════════════════════════════════════
\section{The Direct Hammer: SVD (The ``Small'' Problem)}
\label{sec:svd}

\subsection{When to Use It}

If your matrix fits comfortably in RAM---say $n \le 10{,}000$---the
\textbf{Singular Value Decomposition} is the most robust solver available.

Every $m \times n$ matrix $A$ (real) admits the factorisation
\[
    A = U \Sigma V^T,
\]
where $U \in \mathbb{R}^{m \times m}$ and $V \in \mathbb{R}^{n \times n}$
are orthogonal, and $\Sigma = \operatorname{diag}(\sigma_1, \dots,
\sigma_{\min(m,n)})$ with $\sigma_1 \ge \sigma_2 \ge \cdots \ge 0$.

\subsection{Why SVD?}

\begin{itemize}[nosep]
    \item \textbf{Rank-deficient systems.}  If $A$ is rank-deficient or
          nearly so, LU and Cholesky will amplify noise.  SVD reveals the
          rank through the singular values and silently discards the
          noise subspace.
    \item \textbf{Rectangular systems.}  Least-squares ($m > n$) and
          minimum-norm ($m < n$) solutions fall out of SVD with no extra
          code.
    \item \textbf{Conditioning at a glance.}  The condition number is
          $\kappa(A) = \sigma_1 / \sigma_{\min}$.  If $\kappa \approx
          1/\varepsilon_{\mathrm{mach}}$, you know immediately that
          \emph{no} algorithm can give you accurate digits.
\end{itemize}

\begin{tipbox}
SVD costs $O(\min(m,n)^2 \cdot \max(m,n))$ flops.  For $n = 10{,}000$
that is $\sim 10^{12}$ flops---a few seconds on a modern CPU.
For $n = 100{,}000$ it becomes minutes; for $n = 10^6$ it is
infeasible.  That is when you switch to sparse/iterative methods
(\S\ref{sec:sparse}--\ref{sec:gmres}).
\end{tipbox}

\subsection{The Eigen One-Liner}

\begin{lstlisting}
// Robust solver for any dense matrix
Eigen::VectorXd x = A.bdcSvd(Eigen::ComputeThinU
                            | Eigen::ComputeThinV).solve(b);
\end{lstlisting}

\texttt{bdcSvd} stands for \emph{Bidiagonal Divide-and-Conquer} SVD,
which is Eigen's fastest dense SVD implementation.
\texttt{ComputeThinU/V} tells it to compute only the first
$\min(m,n)$ columns of $U$ and $V$, saving memory and time.

\subsection{Full Example}

\lstinputlisting[caption={Dense SVD solve with rank and condition number reporting.},
                 label={lst:svd}]
                {code/svd_solve.cpp}

% ══════════════════════════════════════════════════════════════════════════
% §5  Sparse Systems & The "Triplet" Rule
% ══════════════════════════════════════════════════════════════════════════
\section{Sparse Systems \& The ``Triplet'' Rule}
\label{sec:sparse}

\subsection{Why Sparse?}

Most matrices that arise from discretising PDEs are \emph{sparse}: the
fraction of non-zero entries is $O(1/n)$ or smaller.  Our heat-rod
tridiagonal has $\sim 3n$ non-zeros out of $n^2$ entries.  A 2-D finite
element mesh on $n$ nodes typically produces $\sim 7n$ non-zeros
(depending on element connectivity).

\begin{center}
\renewcommand{\arraystretch}{1.3}
\begin{tabular}{@{}lcc@{}}
\toprule
& \textbf{Dense} & \textbf{Sparse (CSR)} \\
\midrule
Storage for $10^6 \times 10^6$ & $8\;\text{TB}$ & $\sim 24\;\text{MB}$ \\
Matrix--vector multiply & $O(n^2)$ & $O(\text{nnz})$ \\
Direct solve & $O(n^3)$ & depends on fill-in \\
\bottomrule
\end{tabular}
\end{center}

\subsection{The CSR Format}

\textbf{Compressed Sparse Row (CSR)} stores three arrays:
\begin{itemize}[nosep]
    \item \texttt{values[]}: all non-zero entries, row by row.
    \item \texttt{col\_indices[]}: the column index of each entry.
    \item \texttt{row\_ptr[]}: for each row~$i$, the index into
          \texttt{values} where row~$i$ starts.
\end{itemize}

\begin{example}
For the $3\times 3$ tridiagonal matrix from \S1:
\[
    A = \begin{bmatrix} 2 & -1 & 0 \\ -1 & 2 & -1 \\ 0 & -1 & 2 \end{bmatrix}
\]
\begin{center}
\texttt{values = [2, -1, -1, 2, -1, -1, 2]} \\
\texttt{col\_indices = [0, 1, 0, 1, 2, 1, 2]} \\
\texttt{row\_ptr = [0, 2, 5, 7]}
\end{center}
Total storage: 7 doubles + 7 ints + 4 ints = 100 bytes, versus
72 bytes for the full dense matrix.  The savings become dramatic
when $n > 1000$.
\end{example}

\subsection{The Golden Rule of Assembly}

\begin{warningbox}
\textbf{Never} insert entries one at a time into a compressed sparse matrix:
\begin{lstlisting}
// BAD: O(n^2) total cost -- shifts memory on every insert
for (int i = 0; i < n; ++i)
    A.insert(i, i) = 2.0;
\end{lstlisting}
Each \texttt{insert} may need to shift all subsequent entries to make room.
Over $n$ insertions, this is $O(n^2)$.
\end{warningbox}

\textbf{The Triplet Method:} collect all $(i, j, v)$ entries into a flat
\texttt{std::vector<Triplet>}, then compress once.

\begin{lstlisting}
std::vector<Eigen::Triplet<double>> triplets;
triplets.reserve(3 * n);  // pre-allocate

for (int i = 0; i < n; ++i) {
    triplets.push_back({i, i, 2.0});
    if (i > 0)     triplets.push_back({i, i-1, -1.0});
    if (i < n - 1) triplets.push_back({i, i+1, -1.0});
}

Eigen::SparseMatrix<double> A(n, n);
A.setFromTriplets(triplets.begin(), triplets.end());
\end{lstlisting}

\texttt{setFromTriplets} sorts and compresses in $O(\text{nnz}\log\text{nnz})$.
If duplicate $(i,j)$ pairs exist, their values are \emph{summed} by default---this is exactly what finite element assembly needs.

\subsection{The Triplet Memory Trap}

Each \texttt{Eigen::Triplet<double>} stores three values (row, column,
value)---typically 24~bytes.  For a 1-D rod with $n = 10^6$ and $3n$
non-zeros, that is $\sim\!72$~MB just for the temporary vector.
For a 3-D finite element mesh, the triplet vector can easily exceed the
memory of the final compressed sparse matrix itself.

\begin{tipbox}
Two rules to keep triplet memory under control:
\begin{enumerate}[nosep]
    \item \textbf{Always} call \texttt{triplets.reserve(expected\_nnz)}
          before the assembly loop.  Without it, \texttt{std::vector}
          will re-allocate and copy the entire buffer multiple times,
          temporarily tripling your memory footprint.
    \item \textbf{Free immediately} after compression.  Either scope the
          vector in a block, or call
          \texttt{triplets.clear(); triplets.shrink\_to\_fit();}
          right after \texttt{setFromTriplets}.
\end{enumerate}
\end{tipbox}

\subsection{Full Example: 2-D Poisson on a Grid}

\lstinputlisting[caption={Sparse assembly with triplets for a 2-D Poisson problem.},
                 label={lst:sparse2d}]
                {code/sparse_poisson_2d.cpp}

% ══════════════════════════════════════════════════════════════════════════
% §6  The Iterative Hammer: GMRES
% ══════════════════════════════════════════════════════════════════════════
\section{The Iterative Hammer: GMRES (The ``Big'' Problem)}
\label{sec:gmres}

\subsection{When Direct Solvers Break Down}

Sparse direct solvers (SparseLU, Cholesky) work beautifully up to
$n \sim 10^5$--$10^6$ in 2-D.  In 3-D, however, the factorisation
produces \emph{fill-in}: entries that were zero in $A$ become non-zero
in $L$ and $U$.  For a 3-D mesh with $n$ unknowns, fill-in can grow
as $O(n^{4/3})$, quickly exhausting memory.

Iterative solvers avoid this entirely: they never form $L$ or $U$.
They only need the ability to compute $Ax$ for a given~$x$.

\subsection{How GMRES Works (The 60-Second Version)}

\textbf{Generalised Minimal Residual (GMRES)} builds an orthonormal
basis $\{q_1, q_2, \dots, q_k\}$ for the \emph{Krylov subspace}
\[
    \mathcal{K}_k(A, r_0) = \operatorname{span}\{r_0,\; Ar_0,\;
    A^2 r_0,\; \dots,\; A^{k-1} r_0\},
\]
where $r_0 = b - Ax_0$ is the initial residual.  At step~$k$, it
picks the vector in $x_0 + \mathcal{K}_k$ that minimises the
residual $\|b - Ax\|_2$.

\begin{itemize}[nosep]
    \item \textbf{Convergence:} GMRES converges in at most $n$ steps
          (in exact arithmetic).  In practice it converges much
          sooner if the eigenvalues of $A$ are clustered.
    \item \textbf{Universality:} Unlike Conjugate Gradient (CG),
          GMRES works for \emph{any} non-singular matrix---symmetric
          or not.
    \item \textbf{Memory:} Each iteration stores one new Krylov vector.
          For very large $n$, use \emph{restarted} GMRES($m$),
          which restarts after $m$ vectors.
\end{itemize}

\subsection{Preconditioning: The Secret Sauce}

Without preconditioning, GMRES may need thousands of iterations.
A \textbf{preconditioner} $M \approx A$ transforms the system to
\[
    M^{-1}A\,x = M^{-1}b.
\]
If $M^{-1}A \approx I$, the eigenvalues cluster near~1 and GMRES
converges in very few steps.  The art is choosing $M$ so that:

\begin{enumerate}[nosep]
    \item $M^{-1}v$ is cheap to apply (much cheaper than solving $Ax = v$).
    \item $M^{-1}A$ has clustered eigenvalues.
\end{enumerate}

\begin{definition}[Incomplete LU (ILU)]
Perform LU factorisation of $A$ but \textbf{drop} fill-in entries
below a threshold.  The result $M = \tilde{L}\tilde{U} \approx A$
is cheap to store and apply.  Eigen provides this as
\texttt{IncompleteLUT}.
\end{definition}

\begin{center}
\renewcommand{\arraystretch}{1.3}
\begin{tabular}{@{}lcc@{}}
\toprule
\textbf{Configuration} & \textbf{Typical iterations} & \textbf{Total time} \\
\midrule
No preconditioner      & $1{,}000$--$10{,}000$   & Slow \\
Jacobi (diagonal)      & $500$--$5{,}000$        & Moderate \\
ILU(0)                 & $50$--$200$             & Fast \\
ILU with threshold     & $10$--$50$              & Fastest \\
\bottomrule
\end{tabular}
\end{center}

\begin{warningbox}
A preconditioner that is too aggressive (dropping too many fill-in entries)
can actually \emph{slow down} convergence or cause divergence.
Monitor the residual history; if it stalls or oscillates, tighten the
drop tolerance.
\end{warningbox}

\subsection{GMRES vs.\ BiCGSTAB: Choosing Your Hammer}

Eigen provides two main Krylov solvers.  They are \emph{not}
interchangeable:

\begin{center}
\renewcommand{\arraystretch}{1.3}
\begin{tabular}{@{}lcc@{}}
\toprule
& \textbf{GMRES} & \textbf{BiCGSTAB} \\
\midrule
Robustness        & Indestructible      & Can stall or diverge \\
Memory per iter   & $O(n)$ (stores $q_k$) & $O(n)$ (fixed) \\
Iterations/solve  & Monotone decrease   & Non-monotone (``jittery'') \\
Best for          & Any matrix          & Mildly non-symmetric \\
Failure mode      & Slow, never wrong   & Fast, sometimes explodes \\
\bottomrule
\end{tabular}
\end{center}

\begin{warningbox}
If the matrix is \textbf{highly non-symmetric} or has
\textbf{eigenvalues near the origin}, BiCGSTAB can stagnate or
diverge where GMRES would slowly grind to a correct result.
If you need a ``Magic Hammer'' that never fails, use GMRES.
If you have a deadline and can verify convergence, BiCGSTAB is
often faster per iteration.
\end{warningbox}

\subsection{Eigen Implementation}

Eigen ships \texttt{BiCGSTAB} in the core module and \texttt{GMRES}
in the unsupported module (\texttt{<Eigen/IterativeSolvers>}).
Both accept any preconditioner as a template parameter.  The code
below demonstrates all three configurations---unpreconditioned BiCGSTAB,
ILU-preconditioned BiCGSTAB, and ILU-preconditioned GMRES---on the
same problem so you can compare iteration counts and timing directly.

\lstinputlisting[caption={Iterative solve: BiCGSTAB vs.\ GMRES with ILU preconditioning.},
                 label={lst:gmres}]
                {code/iterative_solve.cpp}

% ══════════════════════════════════════════════════════════════════════════
% §7  The "Don't Do This" Checklist
% ══════════════════════════════════════════════════════════════════════════
\section{The ``Don't Do This'' Checklist}
\label{sec:pitfalls}

Every bug below has crashed a real simulation at least once.
They are ordered from ``most silent'' to ``most spectacular.''

% ── Pitfall 1 ─────────────────────────────────────────────────────────
\subsection{Silent Failure: Ignoring Solver Status}

\begin{warningbox}
Just because a solver function \emph{returns} does not mean it
\emph{succeeded}.  Eigen solvers set an \texttt{info()} flag.
If you never check it, you will silently propagate garbage.
\end{warningbox}

\begin{lstlisting}
// BAD: no error check
Eigen::VectorXd x = solver.solve(b);
use(x);  // x might be nonsense

// GOOD: always check
Eigen::VectorXd x = solver.solve(b);
if (solver.info() != Eigen::Success) {
    std::cerr << "Solver failed: " << solver.info() << "\n";
    return EXIT_FAILURE;
}
\end{lstlisting}

% ── Pitfall 2 ─────────────────────────────────────────────────────────
\subsection{Comparing Floating-Point Values to Zero}

Floating-point arithmetic is \emph{approximate}.  After a chain of
operations, a value that should be exactly zero will instead be
$\sim 10^{-16}$.

\begin{lstlisting}
// BAD: will almost never trigger
if (determinant == 0.0) { /* singular */ }

// ALSO BAD: 1e-12 is an arbitrary "magic number".
// If matrix entries are ~1e-20, then det = 1e-15 is "huge".
// If entries are ~1e+10, then det = 1e-12 is effectively zero.
if (std::abs(determinant) < 1e-12) { /* fragile */ }

// GOOD: use a *relative* tolerance
double scale = A.norm();  // or max entry magnitude
if (std::abs(val) < epsilon * scale) { /* effectively zero */ }
\end{lstlisting}

\begin{tipbox}
Absolute tolerances like \texttt{1e-12} are \textbf{scale-dependent}
and will silently give wrong answers when your matrix entries are
very large or very small.  The most robust check for singularity is the
\emph{condition number} $\kappa(A) = \sigma_1 / \sigma_{\min}$,
which is scale-invariant.
If $\kappa > 1/\varepsilon_{\mathrm{mach}} \approx 10^{16}$,
the matrix is numerically singular regardless of its determinant.
\end{tipbox}

% ── Pitfall 3 ─────────────────────────────────────────────────────────
\subsection{Allocation Inside Hot Loops}

\texttt{Eigen::MatrixXd} dynamically allocates heap memory.
Allocating and freeing inside a tight loop is orders of magnitude
slower than the actual arithmetic.

\begin{lstlisting}
// BAD: allocates n x n on every iteration
while (running) {
    Eigen::MatrixXd A(n, n);   // malloc + free every loop
    fill(A);
    solve(A, b);
}

// GOOD: allocate once, reuse
Eigen::MatrixXd A(n, n);
while (running) {
    fill(A);
    solve(A, b);
}
\end{lstlisting}

For fixed-size small matrices ($n \le 4$), use \texttt{Eigen::Matrix3d}
or \texttt{Eigen::Matrix4d}---they live on the stack and cost zero
allocation overhead.

% ── Pitfall 4 ─────────────────────────────────────────────────────────
\subsection{Skipping the Residual Check}

This is the single most important engineering habit in numerical
computing.  After every solve, compute:
\[
    \text{relative residual} = \frac{\|Ax - b\|}{\|b\|}.
\]

\begin{lstlisting}
Eigen::VectorXd x = solver.solve(b);
double rel_res = (A * x - b).norm() / b.norm();
if (rel_res > 1e-8) {
    std::cerr << "WARNING: residual = " << rel_res
              << " -- solution may be inaccurate\n";
}
\end{lstlisting}

\begin{warningbox}
A small residual does \emph{not} guarantee a small error in~$x$.
If $\kappa(A)$ is large, even a tiny residual can correspond to
a wildly wrong solution.  Always pair the residual check with
a condition number estimate when possible.
\end{warningbox}

% ── Pitfall 5 ─────────────────────────────────────────────────────────
\subsection{Aliasing in Eigen Expressions}

Eigen uses lazy evaluation.  The classic ``silent killer'' is
\texttt{A = A * A}: the same matrix appears on both sides of the
assignment, and without intervention Eigen would read and write
the same memory simultaneously, producing garbage.

In modern Eigen ($\ge 3.3$), the library detects \emph{some}
aliasing patterns automatically---for example,
\texttt{A = A * B} (where \texttt{B} is a \emph{different} matrix)
is handled correctly.  However, \texttt{A = A * A} remains the
classic case that requires explicit intervention.

\begin{lstlisting}
// DANGEROUS: same matrix on both sides
A = A * A;   // Eigen 3.3+ may catch this, but don't rely on it

// SAFE: .eval() is cheap insurance for any self-referencing expression
A = (A * A).eval();
// or equivalently:
Eigen::MatrixXd temp = A * A;
A = temp;
\end{lstlisting}

\begin{tipbox}
Rule of thumb: if the \textbf{same matrix name appears on both sides}
of the equals sign, add \texttt{.eval()}.  The cost is one temporary
allocation---trivial compared to the hours you will spend debugging
silent corruption.
\end{tipbox}

\subsection{Complete Pitfall Demonstration}

\lstinputlisting[caption={All five pitfalls demonstrated and fixed.},
                 label={lst:pitfalls}]
                {code/pitfalls_demo.cpp}


% ── Summary ───────────────────────────────────────────────────────────────
\section{Summary for Busy Engineers}

\begin{center}
\renewcommand{\arraystretch}{1.4}
\begin{tabular}{@{}lll@{}}
\toprule
\textbf{Situation} & \textbf{Tool} & \textbf{Eigen One-Liner} \\
\midrule
Small dense system ($n < 10{,}000$)
    & SVD
    & \texttt{A.bdcSvd(\ldots).solve(b)} \\
Huge sparse system ($n > 100{,}000$)
    & GMRES + ILU
    & \texttt{BiCGSTAB<\ldots,IncompleteLUT>} \\
Slow performance
    & Check memory layout
    & Link to Intel MKL / OpenBLAS \\
\bottomrule
\end{tabular}
\end{center}

\begin{tipbox}
After \emph{every} solve, compute the relative residual
$\|Ax - b\| / \|b\|$. If it is not small, your answer is fiction.
\end{tipbox}

\end{document}
